% !TeX root = ../main.tex
\section{Appendix}
\label{sec:appendix}

\begin{itemize}
    \item Detailed proofs
    \item Maybe more benchmark tables
\end{itemize}

\todo[inline]{Add full benchmark programs here}

\section{Determinization}
\label{sec:appendix-determinization}

We give the full structural definition of the transformation in
\cref{def:determinization}\leanref{determinization}, covering every constructor
in \cref{fig:syntax}. To keep the rules compact, we abbreviate recursive
calls throughout this section by
\[
  e'=\trans{e},
  \qquad e_i'=\trans{e_i}.
\]
Thus, each primed expression denotes the determinization of its unprimed
counterpart. Bound variable names are unchanged.
For each distribution mode $\alpha$, let
\[
  \alpha'=
  \begin{cases}
    \M & \alpha=\E,\\
    \G & \alpha=\G,\\
    \M & \alpha=\M.
  \end{cases}
\]

\paragraph{Variables and constants.}
\begin{align*}
  \trans x &= x,&
  \trans r &= r \quad(r\in\Real),\\
  \trans{()} &= (),&
  \trans{\kw{true}} &= \kw{true},\\
  \trans{\kw{false}} &= \kw{false},&
  \trans{\kw{reject}} &= \kw{reject},\\
  \trans{[]} &= [].
\end{align*}

\paragraph{Functions, recursion, and control flow.}
\begin{align*}
  \trans{\lambda x.e} &= \lambda x.e',\\
  \trans{\recfn f x e} &= \recfn f x{e'},\\
  \trans{e_1\,e_2} &= e_1'\,e_2',\\
  \trans{\letin x{e_1}{e_2}} &= \letin x{e_1'}{e_2'},\\
  \trans{\ite{e_0}{e_1}{e_2}} &= \ite{e_0'}{e_1'}{e_2'}.
\end{align*}

\paragraph{Arithmetic and comparisons.}
\begin{align*}
  \trans{-e} &= -e',\\
  \trans{e_1+e_2} &= e_1'+e_2',\\
  \trans{e_1\cdot e_2} &= e_1'\cdot e_2',\\
  \trans{e_1/e_2} &= e_1'/e_2',\\
  \trans{e_1<e_2} &= e_1'<e_2'.
\end{align*}

\paragraph{Pairs, sums, and lists.}
\begin{align*}
  \trans{\langle e_1,e_2\rangle} &= \langle e_1',e_2'\rangle,\\
  \trans{\kw{fst}\,e} &= \kw{fst}\,e',\\
  \trans{\kw{snd}\,e} &= \kw{snd}\,e',\\
  \trans{\kw{inl}\,e} &= \kw{inl}\,e',\\
  \trans{\kw{inr}\,e} &= \kw{inr}\,e',\\
  \trans{e_1::e_2} &= e_1'::e_2'.
\end{align*}
\begin{align*}
  \trans{\scase{e_0}x{e_1}y{e_2}}
    &={}\\[-2pt]
    &\quad\scase{e_0'}x{e_1'}y{e_2'},\\[4pt]
  \trans{\lcase{e_0}{e_1}x{xs}{e_2}}
    &={}\\[-2pt]
    &\quad\lcase{e_0'}{e_1'}x{xs}{e_2'}.
\end{align*}

\paragraph{Distribution expressions.}
For all $\alpha\in\{\G,\E,\M\}$, the eight primitive distributions transform as
follows:
\begin{align*}
  \trans{\draw{uniform}\alpha{e_1,e_2}}
    &= \draw{uniform}{\alpha'}{e_1',e_2'},\\
  \trans{\draw{gaussian}\alpha{e_1,e_2}}
    &= \draw{gaussian}{\alpha'}{e_1',e_2'},\\
  \trans{\draw{poisson}\alpha e}
    &= \draw{poisson}{\alpha'}{e'},\\
  \trans{\draw{exponential}\alpha e}
    &= \draw{exponential}{\alpha'}{e'},\\
  \trans{\draw{bernoulli}\alpha e}
    &= \draw{bernoulli}{\alpha'}{e'},\\
  \trans{\draw{discrete}\alpha e}
    &= \draw{discrete}{\alpha'}{e'},\\
  \trans{\draw{beta}\alpha{e_1,e_2}}
    &= \draw{beta}{\alpha'}{e_1',e_2'},\\
  \trans{\draw{gamma}\alpha{e_1,e_2}}
    &= \draw{gamma}{\alpha'}{e_1',e_2'}.
\end{align*}
The target retains each distribution primitive and all its arguments; only
$\E$ changes to $\M$, and every argument is recursively transformed.
Mean evaluation remains an atomic distribution operation, governed by the
domain conditions and reduction rule in \cref{fig:semantics}.

\section{Miscellaneous}\label{sec:appendix-misc}
\paragraph{When Expectations Are Undefined}
\begin{example}
In the coin-flipping example (cf. \Cref{ex:coin-flip}), each terminating $\G$-trace has positive probability. A trace can also contain continuous draws:
\begin{lstlisting}
let u = uniform[G](-1, 1) in
uniform[E](0, 2) / u
\end{lstlisting}
This program has an undefined global expectation. Yet for each fixed $u\ne0$,
the remaining expression has mean $1/u$, exactly the value returned by the
transformed program. Thus determinization is tracewise sound for this program, but the global expectation preservation theorem does not apply. 

Moreover, note that every individual value of $u$ has probability zero, so conditioning on a trace cannot simply mean dividing by its probability, hence the disintegration-based definition of conditional output distributions in~\Cref{eq:disintegration}.
\end{example}

\begin{example}
Consider
\begin{lstlisting}
let u = uniform[G](0, 1) in
let s = 2 * bernoulli[G](1/2) - 1 in
s / u + uniform[E](-1, 1)
\end{lstlisting}
The program returns almost surely and is domain-safe: although division is
undefined when $u=0$, this event has probability zero under the
continuous uniform distribution. However,
\[
  \int_0^1 \frac{1}{u}\,du=\infty.
\]
When $s=1$, the term $s/u$ gives the return value an infinite positive part;
when $s=-1$, it gives the return value an infinite negative part. Since each
sign occurs with probability $1/2$, both parts have infinite expectation, so
the program's global expectation is undefined.

Determinization is nevertheless valid and tracewise sound for this program. Determinization replaces the final uniform draw by its mean zero and therefore
returns exactly $s/u$ for some fixed $\G$-trace. The global expectation preservation theorem does not
apply because its hypothesis requires the source expectation to be
well-defined in the extended reals. Likewise, the variance theorem does not
apply because the program has an infinite second moment.
\end{example}

\paragraph{Domain safety and missing probability.}
\begin{example}
The following program is not domain-safe:
\begin{lstlisting}
let x = bernoulli[G](1/2) in
1 / x
\end{lstlisting}
With probability $1/2$, the sample is $x=0$, at which point division by zero
causes execution to get stuck. Because this invalid operation occurs with
positive probability, it cannot be ignored by the almost-sure condition in
domain safety. When $x=1$, the program successfully returns $1$, so
\[
  \law e=\frac12\delta_1,
  \qquad
  \law e(\Real)=\frac12.
\]
The missing probability mass comes from an invalid operation rather than
rejection or divergence.
\end{example}

\begin{example}
The following well-typed program $e$ is not domain-safe:
\begin{lstlisting}
let u = gaussian[E](0, 1) in uniform[E](0, u)
\end{lstlisting}
When $u<0$, the uniform distribution's parameters are invalid and execution gets
stuck. The program returns with probability $1/2$; conditioned on returning,
its expected output is half the mean of a standard Gaussian conditioned to be
positive. Determinization sets $u=0$ and returns the mean of
$\kw{uniform}(0,0)$, which is zero.
Thus the domain-safety assumption is necessary, because:
\[
 \Expect_{\mathrm{ret}}[e]=\frac{1}{\sqrt{2\pi}}
 \qquad\text{but}\qquad
 \Expect_{\mathrm{ret}}[\trans e]=0.
\]
\end{example}

\paragraph{Why domain safety is necessary.}
The preceding guarantees rely essentially on domain safety. Without it, replacing an $\E$-mode sample by its mean can change whether a later operation is defined, and therefore change which executions contribute to the return-conditioned distribution.

For example,
\begin{lstlisting}
let u = gaussian[E](0, 1) in
uniform[E](0, u)
\end{lstlisting}
is not domain-safe: executions with $u<0$ become stuck. Conditioning on successful return therefore implicitly conditions $u$ to be positive, giving
\[
\Expect_{\mathrm{ret}}[e]
=
\frac{1}{2}\Expect[u\mid u>0]
=
\frac{1}{\sqrt{2\pi}}.
\]
Determinization instead replaces $u$ by its unconditional mean $0$, yielding
\[
\Expect_{\mathrm{ret}}[\trans e]=0.
\]
Thus domain safety prevents determinization from changing which executions are semantically valid and, consequently, which executions contribute to the return-conditioned output distribution.


\subsection{Formal Definitions for Symbolic Execution}
\label{sec:appendix-symbolic-definitions}

\begin{definition}[Residual symbolic expressions and states]
Fix symbolic variables $u_1,\ldots,u_k$. An affine expression, a residual
symbolic expression, a symbolic environment, and a symbolic state are given by
the following grammar:
\[
\begin{aligned}
  a
  &::=
  c_0+\sum_{i=1}^{k}c_i u_i,
  \qquad c_0,\ldots,c_k\in\Real,
  \\
  \widehat e
  &::=
  x
  \mid a
  \mid ()
  \mid \kw{true}
  \mid \kw{false}
  \mid \kw{reject}
  \\
  &\quad\mid
  \lambda x.\widehat e
  \mid \recfn f x{\widehat e}
  \mid \widehat e\,\widehat e
  \mid \letin x{\widehat e}{\widehat e}
  \\
  &\quad\mid
  \ite{\widehat e}{\widehat e}{\widehat e}
  \mid -\widehat e
  \mid \widehat e+\widehat e
  \mid \widehat e\cdot\widehat e
  \mid \widehat e/\widehat e
  \mid \widehat e<\widehat e
  \\
  &\quad\mid
  \langle\widehat e,\widehat e\rangle
  \mid \kw{fst}\,\widehat e
  \mid \kw{snd}\,\widehat e
  \\
  &\quad\mid
  \kw{inl}\,\widehat e
  \mid \kw{inr}\,\widehat e
  \mid \scase{\widehat e}{x}{\widehat e_1}{y}{\widehat e_2}
  \\
  &\quad\mid
  []
  \mid \widehat e::\widehat e
  \mid \lcase{\widehat e}{\widehat e_0}{x}{xs}{\widehat e_1}
  \\
  &\quad\mid
  \draw{uniform}\alpha{\widehat e,\widehat e}
  \mid \draw{gaussian}\alpha{\widehat e,\widehat e}
  \mid \draw{poisson}\alpha{\widehat e}
  \mid \draw{exponential}\alpha{\widehat e}
  \\
  &\quad\mid
  \draw{bernoulli}\alpha{\widehat e}
  \mid \draw{discrete}\alpha{\widehat e}
  \mid \draw{beta}\alpha{\widehat e,\widehat e}
  \mid \draw{gamma}\alpha{\widehat e,\widehat e},
  \\
  \alpha
  &::=\G\mid\E\mid\M,
  \\
  \sigma
  &::=
  \emptyset
  \mid
  \sigma,\ u\sim D(\bar a),
  \\
  s
  &::=
  (\sigma\mid\widehat e:\tau).
\end{aligned}
\]
Thus, the residual grammar is exactly the source grammar with each real
literal replaced by an affine expression.  In a state containing $k$
declarations, every affine expression has precisely the coordinates
$u_1,\ldots,u_k$.  A declaration $u_i\sim D_i(\bar a_i)$ is well formed only
when $u_i$ is fresh and every parameter in $\bar a_i$ refers only to
$u_1,\ldots,u_{i-1}$.  When symbolic execution appends a declaration, the
affine expressions already present in the state are extended with coefficient
zero for the fresh coordinate.  Any distribution parameter required to have
type $\tyR\G$ is a constant affine expression by the typing rules below.
\end{definition}

\begin{definition}[Typing residual symbolic expressions]
Residual symbolic expressions use the type grammar from \cref{fig:typing}.
The judgment $\typing{\widehat e}\tau$ is generated by the following rules.
The two affine-literal rules are the only rules that differ from the source
typing judgment.

\begin{mathpar}
  \inferrule*[right=Affine-E]
    {a=c_0+\sum_{i=1}^{k}c_i u_i}
    {\typing a{\tyR\E}}
  \and
  \inferrule*[right=Affine-G]
    {a=c_0+\sum_{i=1}^{k}c_i u_i\\c_1=\cdots=c_k=0}
    {\typing a{\tyR\G}}
  \and
  \inferrule*[right=Var]
    {x:\tau\in\Gamma}
    {\typing x\tau}
  \and
  \inferrule*[right=Unit]
    { }
    {\typing{()}\tyU}
  \and
  \inferrule*[right=True]
    { }
    {\typing{\kw{true}}\tyB}
  \and
  \inferrule*[right=False]
    { }
    {\typing{\kw{false}}\tyB}
  \and
  \inferrule*[right=Reject]
    { }
    {\typing{\kw{reject}}\tau}
\end{mathpar}

\begin{mathpar}
  \inferrule*[right=Lam]
    {\Gamma,x:\tau_1\vdash\widehat e:\tau_2}
    {\typing{\lambda x.\widehat e}{\tau_1\to\tau_2}}
  \and
  \inferrule*[right=Rec]
    {\Gamma,f:\tau_1\to\tau_2,x:\tau_1\vdash\widehat e:\tau_2}
    {\typing{\recfn f x{\widehat e}}{\tau_1\to\tau_2}}
  \and
  \inferrule*[right=App]
    {\typing{\widehat e_1}{\tau_1\to\tau_2}\\
     \typing{\widehat e_2}{\tau_1}}
    {\typing{\widehat e_1\,\widehat e_2}{\tau_2}}
  \and
  \inferrule*[right=Let]
    {\typing{\widehat e_1}{\tau_1}\\
     \Gamma,x:\tau_1\vdash\widehat e_2:\tau_2}
    {\typing{\letin x{\widehat e_1}{\widehat e_2}}{\tau_2}}
  \and
  \inferrule*[right=If]
    {\typing{\widehat e_1}\tyB\\
     \typing{\widehat e_2}\tau\\
     \typing{\widehat e_3}\tau}
    {\typing{\ite{\widehat e_1}{\widehat e_2}{\widehat e_3}}\tau}
\end{mathpar}

\begin{mathpar}
  \inferrule*[right=Pair]
    {\typing{\widehat e_1}{\tau_1}\\
     \typing{\widehat e_2}{\tau_2}}
    {\typing{\langle\widehat e_1,\widehat e_2\rangle}{\tau_1\times\tau_2}}
  \and
  \inferrule*[right=Fst]
    {\typing{\widehat e}{\tau_1\times\tau_2}}
    {\typing{\kw{fst}\,\widehat e}{\tau_1}}
  \and
  \inferrule*[right=Snd]
    {\typing{\widehat e}{\tau_1\times\tau_2}}
    {\typing{\kw{snd}\,\widehat e}{\tau_2}}
  \and
  \inferrule*[right=Inl]
    {\typing{\widehat e}{\tau_1}}
    {\typing{\kw{inl}\,\widehat e}{\tau_1+\tau_2}}
  \and
  \inferrule*[right=Inr]
    {\typing{\widehat e}{\tau_2}}
    {\typing{\kw{inr}\,\widehat e}{\tau_1+\tau_2}}
  \and
  \inferrule*[right=Sum-Case]
    {\typing{\widehat e}{\tau_1+\tau_2}\\
     \Gamma,x:\tau_1\vdash\widehat e_1:\tau\\
     \Gamma,y:\tau_2\vdash\widehat e_2:\tau}
    {\typing{\scase{\widehat e}{x}{\widehat e_1}{y}{\widehat e_2}}\tau}
\end{mathpar}

\begin{mathpar}
  \inferrule*[right=Nil]
    { }
    {\typing{[]}{\tyL\tau}}
  \and
  \inferrule*[right=Cons]
    {\typing{\widehat e_1}\tau\\
     \typing{\widehat e_2}{\tyL\tau}}
    {\typing{\widehat e_1::\widehat e_2}{\tyL\tau}}
  \and
  \inferrule*[right=List-Case]
    {\typing{\widehat e}{\tyL{\tau_1}}\\
     \typing{\widehat e_0}{\tau_2}\\
     \Gamma,x:\tau_1,xs:\tyL{\tau_1}\vdash\widehat e_1:\tau_2}
    {\typing{\lcase{\widehat e}{\widehat e_0}{x}{xs}{\widehat e_1}}{\tau_2}}
\end{mathpar}

\begin{mathpar}
  \inferrule*[right=Neg]
    {\typing{\widehat e}{\tyR m}}
    {\typing{-\widehat e}{\tyR m}}
  \and
  \inferrule*[right=Add]
    {\typing{\widehat e_1}{\tyR m}\\
     \typing{\widehat e_2}{\tyR m}}
    {\typing{\widehat e_1+\widehat e_2}{\tyR m}}
  \and
  \inferrule*[right=Mul]
    {\typing{\widehat e_1}{\tyR\G}\\
     \typing{\widehat e_2}{\tyR m}}
    {\typing{\widehat e_1\cdot\widehat e_2}{\tyR m}}
  \and
  \inferrule*[right=Div]
    {\typing{\widehat e_1}{\tyR m}\\
     \typing{\widehat e_2}{\tyR\G}}
    {\typing{\widehat e_1/\widehat e_2}{\tyR m}}
  \and
  \inferrule*[right=Cmp]
    {\typing{\widehat e_1}{\tyR\G}\\
     \typing{\widehat e_2}{\tyR\G}}
    {\typing{\widehat e_1<\widehat e_2}\tyB}
\end{mathpar}

For the distribution rules below, $m\in\{\G,\E\}$ and
$\alpha\in\{m,\M\}$.
\begin{mathpar}
  \inferrule*[right=Uniform]
    {\typing{\widehat e_1}{\tyR m}\\
     \typing{\widehat e_2}{\tyR m}}
    {\typing{\draw{uniform}\alpha{\widehat e_1,\widehat e_2}}{\tyR m}}
  \and
  \inferrule*[right=Gaussian]
    {\typing{\widehat e_1}{\tyR m}\\
     \typing{\widehat e_2}{\tyR\G}}
    {\typing{\draw{gaussian}\alpha{\widehat e_1,\widehat e_2}}{\tyR m}}
  \and
  \inferrule*[right=Poisson]
    {\typing{\widehat e}{\tyR m}}
    {\typing{\draw{poisson}\alpha{\widehat e}}{\tyR m}}
  \and
  \inferrule*[right=Exponential]
    {\typing{\widehat e}{\tyR\G}}
    {\typing{\draw{exponential}\alpha{\widehat e}}{\tyR m}}
  \and
  \inferrule*[right=Bernoulli]
    {\typing{\widehat e}{\tyR m}}
    {\typing{\draw{bernoulli}\alpha{\widehat e}}{\tyR m}}
  \and
  \inferrule*[right=Discrete]
    {\typing{\widehat e}{\tyL{\tyR m}}}
    {\typing{\draw{discrete}\alpha{\widehat e}}{\tyR m}}
  \and
  \inferrule*[right=Beta]
    {\typing{\widehat e_1}{\tyR\G}\\
     \typing{\widehat e_2}{\tyR\G}}
    {\typing{\draw{beta}\alpha{\widehat e_1,\widehat e_2}}{\tyR m}}
  \and
  \inferrule*[right=Gamma]
    {\typing{\widehat e_1}{\tyR m}\\
     \typing{\widehat e_2}{\tyR\G}}
    {\typing{\draw{gamma}\alpha{\widehat e_1,\widehat e_2}}{\tyR m}}
\end{mathpar}

The subtyping relation is unchanged:
\begin{mathpar}
  \inferrule*[right=Refl]
    { }
    {\tau<:\tau}
  \and
  \inferrule*[right=Trans]
    {\tau<:\sigma\\\sigma<:\rho}
    {\tau<:\rho}
  \and
  \inferrule*[right=Real-Sub]
    { }
    {\tyR\G<:\tyR\E}
  \and
  \inferrule*[right=Arrow]
    {\sigma_1<:\tau_1\\\tau_2<:\sigma_2}
    {\tau_1\to\tau_2<:\sigma_1\to\sigma_2}
  \and
  \inferrule*[right=Product]
    {\tau_1<:\sigma_1\\\tau_2<:\sigma_2}
    {\tau_1\times\tau_2<:\sigma_1\times\sigma_2}
  \and
  \inferrule*[right=Sum]
    {\tau_1<:\sigma_1\\\tau_2<:\sigma_2}
    {\tau_1+\tau_2<:\sigma_1+\sigma_2}
  \and
  \inferrule*[right=List]
    {\tau<:\sigma}
    {\tyL\tau<:\tyL\sigma}
  \and
  \inferrule*[right=Sub]
    {\typing{\widehat e}\tau\\\tau<:\sigma}
    {\typing{\widehat e}\sigma}
\end{mathpar}
A symbolic state $(\sigma\mid\widehat e:\tau)$ is well typed when
$\emptyset\vdash\widehat e:\tau$ under these rules.  The declarations in
$\sigma$ supply the coordinates of affine literals; they are not term
variables in $\Gamma$.
\end{definition}


\begin{definition}[Lifting, realization, and initialization]
The lifting function $\mathsf{lift}$ maps an ordinary expression to a
residual symbolic expression.  It maps a real literal $r$ to the constant
affine expression with constant coefficient $r$ and no variable
coefficients, and it acts recursively on every other constructor.  In
particular,
\[
\begin{aligned}
  \mathsf{lift}(r)
  &=r,
  \\
  \mathsf{lift}(e_1+e_2)
  &=\mathsf{lift}(e_1)+\mathsf{lift}(e_2),
  \\
  \mathsf{lift}\bigl(D^\alpha(e_1,\ldots,e_n)\bigr)
  &=D^\alpha\bigl(
      \mathsf{lift}(e_1),\ldots,\mathsf{lift}(e_n)
    \bigr).
\end{aligned}
\]
The first $r$ above is an ordinary literal, while the second is its constant
affine representation.

For a state with $k$ declarations, a valuation $\rho$ assigns a real value
$\rho(u_i)$ to each $u_i$.  Realization maps a residual symbolic expression
back to an ordinary expression by
\[
\begin{aligned}
  \mathsf{realize}_\rho
  \left(c_0+\sum_{i=1}^{k}c_i u_i\right)
  &=c_0+\sum_{i=1}^{k}c_i\rho(u_i),
  \\
  \mathsf{realize}_\rho(\widehat e_1+\widehat e_2)
  &=\mathsf{realize}_\rho(\widehat e_1)
    +\mathsf{realize}_\rho(\widehat e_2),
  \\
  \mathsf{realize}_\rho
  \bigl(D^\alpha(\widehat e_1,\ldots,\widehat e_n)\bigr)
  &=D^\alpha\bigl(
      \mathsf{realize}_\rho(\widehat e_1),\ldots,
      \mathsf{realize}_\rho(\widehat e_n)
    \bigr).
\end{aligned}
\]
It acts recursively in the same way on every remaining constructor, preserving
binders and distribution modes.  Finally,
\[
  \mathsf{init}(e)
  :=
  \bigl(\emptyset\mid\mathsf{lift}(e):\tau\bigr)
  \qquad\text{when }\emptyset\vdash e:\tau.
\]
\end{definition}

\begin{definition}[Actual environment law and sequential mean]
Let
\[
  \sigma
  =
  \bigl(
    u_1\sim D_1(\bar a_1),\ldots,
    u_k\sim D_k(\bar a_k)
  \bigr).
\]
The law $\mathcal A_\sigma$ of its valuation is defined recursively.  For
every measurable set $B$ of valuations,
\[
\begin{aligned}
  \mathcal A_\emptyset
  &:=\delta_{()},
  \\
  \mathcal A_{\sigma,u\sim D(\bar a)}(B)
  &:=
  \int
    D\bigl(\bar a[\theta]\bigr)
      \bigl(\{r\in\Real\mid(\theta,r)\in B\}\bigr)\,
    \mathcal A_\sigma(d\theta).
\end{aligned}
\]
Here $(\theta,r)$ denotes the valuation that extends $\theta$ by assigning
$r$ to the fresh variable $u$.
Thus a valuation is sampled from $\mathcal A_\sigma$ before the parameters
$\bar a[\theta]$ of the next declaration are evaluated.  Domain safety
ensures that each distribution in this recursion is defined and that
$\mathcal A_\sigma$ is a probability measure.

The sequential mean valuation is defined in the same declaration order:
\[
\begin{aligned}
  \bar\theta_\emptyset
  &:=(),
  \\
  \bar\theta_{\sigma,u\sim D(\bar a)}
  &:=
  \bigl(
    \bar\theta_\sigma,
    \operatorname{Mean}_D
      (\bar a[\bar\theta_\sigma])
  \bigr).
\end{aligned}
\]
Consequently,
\[
  \bar\theta_\sigma(u_i)
  =
  \operatorname{Mean}_{D_i}
  \bigl(\bar a_i[\bar\theta_\sigma]\bigr),
  \qquad i=1,\ldots,k,
\]
where the right-hand side depends only on the values already assigned to
$u_1,\ldots,u_{i-1}$.
\end{definition}

\begin{definition}[Actual and expected interpretations]
For a symbolic state $s=(\sigma\mid\widehat e:\tau)$, define
\[
\begin{aligned}
  \mathsf{Act}_n(s)
  &:={}
  \int
    \joint{\mathsf{realize}_\theta(\widehat e)}n\,
    \mathcal A_\sigma(d\theta),
  \\
  \mathsf{Exp}_n(s)
  &:={}
  \joint{
    \trans{
      \mathsf{realize}_{\bar\theta_\sigma}(\widehat e)
    }
  }n.
\end{aligned}
\]
Both are measures on $\Trace\times\Real$.  The actual interpretation samples
the delayed $\E$-mode declarations according to $\mathcal A_\sigma$; the
expected interpretation realizes them at their sequential means and then
determinizes any $\E$-mode sites that symbolic execution has not yet reached.
\end{definition}

\subsection{Theorems and Lemmas from the Symbolic Semantics Proof}

\begin{lemma}[Lifting preserves typing]
\label{lem:appendix-lift-typing}
If $\typing e\tau$, then
\[
  \typing{\mathsf{lift}(e)}\tau.
\]
\end{lemma}

\begin{lemma}[Realization preserves typing]
\label{lem:appendix-realize-typing}
If $\typing{\widehat e}\tau$, then, for every valuation $\rho$,
\[
  \typing{\mathsf{realize}_\rho(\widehat e)}\tau.
\]
\end{lemma}

\begin{lemma}[Determinization preserves typing]
\label{lem:appendix-determinization-typing}
If $\typing e\tau$, then
\[
  \typing{\trans e}\tau.
\]
\end{lemma}

\begin{lemma}[Realization of initialization]
\label{lem:appendix-realize-initialization}
For every ordinary expression $e$,
\[
  \mathsf{realize}_{()}\bigl(\mathsf{lift}(e)\bigr)=e.
\]
Consequently, if $\emptyset\vdash e:\tau$, then
$\mathsf{init}(e)$ is a well-typed symbolic state of type $\tau$.
\end{lemma}

\begin{lemma}[The actual environment law is a probability measure]
\label{lem:appendix-actual-environment-law}
If $\sigma$ is domain-safe, then $\mathcal A_\sigma$ is a probability
measure on valuations of the variables declared in $\sigma$.
\end{lemma}

\begin{lemma}[Affine means]
\label{lem:appendix-affine-means}
Suppose $\sigma$ is domain-safe and $a$ is an affine expression over the
variables declared in $\sigma$. Then
\[
  \int |a[\theta]|\,\mathcal A_\sigma(d\theta)<\infty,
  \qquad
  \int a[\theta]\,\mathcal A_\sigma(d\theta)
  =
  a[\bar\theta_\sigma].
\]
\end{lemma}

\begin{lemma}[Preservation by symbolic reduction]
\label{lem:appendix-symbolic-preservation}
Suppose $s$ is well typed and domain-safe. A symbolic state produced from
$s$ by a \textsc{Reduce} or \textsc{Sample-E} step is also well typed and
domain-safe. For a \textsc{Sample-G} step, the resulting state is well typed
and domain-safe for $D(\bar r)$-almost every sampled value $r$.
\end{lemma}

\begin{lemma}[Actual interpretation commutes with a symbolic step]
\label{lem:appendix-actual-commuting}
Let $s$ be a well-typed, domain-safe symbolic state. If a \textsc{Reduce} or
\textsc{Sample-E} step gives
\[
  s\xrightarrow{}_{\mathrm{sym}}s',
\]
then, for every $n\in\Nat$,
\[
  \mathsf{Act}_{n+1}(s)=\mathsf{Act}_n(s').
\]
For a \textsc{Sample-G} state
\[
  s=(\sigma\mid\widehat C[\draw D\G{\bar r}]:\tau),
\]
every measurable $A\subseteq\Trace$ and $B\subseteq\Real$ satisfy
\[
\begin{aligned}
  \mathsf{Act}_{n+1}(s)(A\times B)
  ={}&
  \int_{\Real}
  \mathsf{Act}_n
    \bigl((\sigma\mid\widehat C[r]:\tau)\bigr)
    \bigl(
      \{(\tau',x)\mid(D,r)::\tau'\in A\ \text{and}\ x\in B\}
    \bigr)
  \,D(\bar r)(dr).
\end{aligned}
\]
\end{lemma}

\begin{lemma}[Expected interpretation commutes with a symbolic step]
\label{lem:appendix-expected-commuting}
Let $s$ be a well-typed, domain-safe symbolic state. If a \textsc{Reduce} or
\textsc{Sample-E} step gives
\[
  s\xrightarrow{}_{\mathrm{sym}}s',
\]
then, for every $n\in\Nat$,
\[
  \mathsf{Exp}_{n+1}(s)=\mathsf{Exp}_n(s').
\]
For a \textsc{Sample-G} state
\[
  s=(\sigma\mid\widehat C[\draw D\G{\bar r}]:\tau),
\]
every measurable $A\subseteq\Trace$ and $B\subseteq\Real$ satisfy
\[
\begin{aligned}
  \mathsf{Exp}_{n+1}(s)(A\times B)
  ={}&
  \int_{\Real}
  \mathsf{Exp}_n
    \bigl((\sigma\mid\widehat C[r]:\tau)\bigr)
    \bigl(
      \{(\tau',x)\mid(D,r)::\tau'\in A\ \text{and}\ x\in B\}
    \bigr)
  \,D(\bar r)(dr).
\end{aligned}
\]
\end{lemma}

\begin{theorem}[Interpretation of initialization]
\label{thm:appendix-interpretation-initialization}
Suppose $\vdash e:\tyR\E$ and $e$ is closed. Then, for every $n\in\Nat$,
\[
  \mathsf{Act}_n\bigl(\mathsf{init}(e)\bigr)=\joint e n,
  \qquad
  \mathsf{Exp}_n\bigl(\mathsf{init}(e)\bigr)=\joint{\trans e}n.
\]
\end{theorem}

\begin{theorem}[Agreement of the two interpretations]
\label{thm:appendix-interpretation-agreement}
Let $s$ be a well-typed, domain-safe symbolic state. For every $n\in\Nat$
and every measurable $A\subseteq\Trace$,
\[
  \mathsf{Act}_n(s)(A\times\Real)
  =
  \mathsf{Exp}_n(s)(A\times\Real).
\]
Moreover, for almost every trace $\tau$ under this common trace marginal,
\[
  \int_{\Real}|r|\,\mathsf{Act}_n(s)(dr\mid\tau)<\infty,
  \qquad
  \mathsf{Exp}_n(s)(\,\cdot\mid\tau)
  =
  \delta_{\displaystyle
    \int_{\Real}r\,\mathsf{Act}_n(s)(dr\mid\tau)}.
\]
\end{theorem}

\begin{theorem}[Agreement after $n$ steps]
\label{thm:appendix-initialized-agreement}
Suppose $\vdash e:\tyR\E$ and $e$ is closed and domain-safe. For every
$n\in\Nat$ and every measurable $A\subseteq\Trace$,
\[
  \joint e n(A\times\Real)
  =
  \joint{\trans e}n(A\times\Real).
\]
Moreover, for almost every trace $\tau$ under this common trace marginal,
\[
  \int_{\Real}|r|\,\joint e n(dr\mid\tau)<\infty,
  \qquad
  \joint{\trans e}n(\,\cdot\mid\tau)
  =
  \delta_{\displaystyle
    \int_{\Real}r\,\joint e n(dr\mid\tau)}.
\]
\end{theorem}

\begin{lemma}[Big-step interpretation of initialization]
\label{lem:appendix-big-step-initialization}
Suppose $\vdash e:\tyR\E$ and $e$ is closed and domain-safe. Then
\[
  \sum_{n=0}^{\infty}
    \mathsf{Act}_n\bigl(\mathsf{init}(e)\bigr)
  =J_e,
  \qquad
  \sum_{n=0}^{\infty}
    \mathsf{Exp}_n\bigl(\mathsf{init}(e)\bigr)
  =J_{\trans e}.
\]
\end{lemma}

\begin{lemma}[Big-step symbolic soundness]
\label{lem:appendix-big-step-symbolic-soundness}
Suppose $\vdash e:\tyR\E$ and $e$ is closed and domain-safe. Then, for
$\trlaw e$-almost every trace $\tau$,
\[
  \int_{\Real}|r|\,\law e(dr\mid\tau)<\infty.
\]
Moreover, for all measurable $A\subseteq\Trace$ and
$B\subseteq\Real$,
\[
  J_{\trans e}(A\times B)
  =
  \int_A
    \delta_{\Expect[e\mid\tau]}(B)\,
    \trlaw e(d\tau).
\]
\end{lemma}

\begin{lemma}[Trace preservation]
Suppose $\vdash e:\tyR\E$ and $e$ is domain-safe. Then $\trans e$ is
domain-safe and
\[
  \trlaw{\trans e}=\trlaw e.
\]
\end{lemma}

\begin{theorem}[Tracewise soundness]
Suppose $\vdash e:\tyR\E$ and $e$ is domain-safe. For
$\trlaw e$-almost every trace $\tau$,
\[
  \law{\trans e}(\,\cdot\mid\tau)
  =
  \delta_{\Expect[e\mid\tau]}.
\]
In particular,
\[
  \Expect[e\mid\tau],\ \Expect[\trans e\mid\tau]\in\Real,
  \qquad
  \Expect[\trans e\mid\tau]=\Expect[e\mid\tau].
\]
\end{theorem}

\subsection{Global Soundness}

\begin{lemma}[Output-mass preservation]
Suppose $\vdash e:\tyR\E$ and $e$ is domain-safe. Then
\[
  q_{\trans e}=q_e.
\]
\end{lemma}

\begin{theorem}[Expectation preservation]
Suppose $\vdash e:\tyR\E$ and $e$ is domain-safe. If $q_e>0$ and
$\Expect_{\mathrm{ret}}[e]$ is well-defined in the extended reals, then
$\Expect_{\mathrm{ret}}[\trans e]$ is also well-defined and
\[
  \Expect_{\mathrm{ret}}[\trans e]
  =
  \Expect_{\mathrm{ret}}[e].
\]
\end{theorem}

\begin{theorem}[Convex function inequality]
Suppose $\vdash e:\tyR\E$ and $e$ is domain-safe. For every nonnegative
convex function $\varphi:\Real\to\Real$,
\[
  \int_{\Real}\varphi(r)\,\law{\trans e}(dr)
  \leq
  \int_{\Real}\varphi(r)\,\law e(dr),
\]
where both sides are nonnegative extended integrals and may equal
$+\infty$.
\end{theorem}

\paragraph{Variance decomposition.}
The transformed program retains only variation between the conditional means.
The original program's output can also vary after the $\G$ trace is fixed,
due to the $\E$-samples.

\begin{theorem}[Trace variance decomposition\leanref{trace-variance}]
\label{thm:trace-variance}
Suppose $\vdash e:\tyR\E$ and $e$ is domain-safe. If
$q_e=\law e(\Real)>0$ and $\int_{\Real}r^2\,\law e(dr)<\infty$, then
\[
 \underbrace{\Var_{\mathrm{ret}}[e]}_{\text{Original variance}}
 =\underbrace{\Var_{\mathrm{ret}}[\trans e]}_{\text{Variance between traces}}
  +\underbrace{\Expect_{\tau\sim\trlaw e/q_e}
    \!\left[\Var\bigl(\law e(\,\cdot\mid\tau)\bigr)\right]}_{\text{Average variance within a trace}}.
\]
\end{theorem}

The last term averages the original program's variance after fixing a trace, using
$\trlaw e/q_e$ to condition on returning. It is nonnegative and gives exactly
the variance removed by determinization.

\begin{theorem}[Variance non-increase]
Suppose $\vdash e:\tyR\E$ and $e$ is domain-safe. If $q_e>0$ and
\[
  \int_{\Real}r^2\,\law e(dr)<\infty,
\]
then $\trans e$ has finite second moment and
\[
  \Var_{\mathrm{ret}}[\trans e]
  \leq
  \Var_{\mathrm{ret}}[e].
\]
\end{theorem}

\begin{proposition}[Probability preservation]
Suppose $\vdash e:\tyR\E$ and $e$ is domain-safe. Then $\trans e$ is
domain-safe and
\[
  \law e(\Real)+d_e=1,
  \qquad
  \law{\trans e}(\Real)=\law e(\Real),
  \qquad
  d_{\trans e}=d_e.
\]
Consequently, determinization preserves the combined probability of
rejection and divergence.
\end{proposition}
